%%=============================================================================
%% 附录 A 实验超参数与配置
%%=============================================================================

\xmuchapter{附录}{Appendix}
\label{chap:appendix}

本附录详细列出本研究中所有实验的超参数设置与软件环境配置，以确保研究的可复现性。所有超参数均在统一的配置文件中管理，便于调试和修改。

\xmusection{数据参数与数据集}{Data Parameters and Dataset}
\label{sec:data_and_dataset}

EEG 数据采集与预处理的基础参数如表~\ref{tab:eeg_params} 所示。

\begin{table}[H]
\centering
\caption{EEG 数据参数设置}
\label{tab:eeg_params}
\begin{tabular}{ll}
\hline
\textbf{参数} & \textbf{值} \\
\hline
采样频率 & 250 Hz \\
通道数 & 22 \\
类别数 & 4（左手、右手、脚、舌头） \\
试次持续时间 & 4.0 秒 \\
带通滤波范围 & 8--30 Hz \\
陷波滤波频率 & 50 Hz \\
\hline
\end{tabular}
\end{table}

本研究使用的数据集信息如表~\ref{tab:dataset_info} 所示。

\begin{table}[H]
\centering
\caption{BCI Competition IV 2a 数据集信息}
\label{tab:dataset_info}
\begin{tabular}{ll}
\hline
\textbf{属性} & \textbf{描述} \\
\hline
数据集名称 & BCI Competition IV 2a \\
被试数量 & 9 名健康被试 \\
任务类型 & 4 类运动想象（左手、右手、脚、舌头） \\
EEG 通道数 & 22 通道 \\
采样频率 & 250 Hz \\
通道位置 & 国际标准 10-20 系统 \\
训练集试次数 & 288 试次/被试（72 试次 × 4 类） \\
测试集试次数 & 288 试次/被试（72 试次 × 4 类） \\
数据来源 & \url{https://github.com/bregydoc/bcidatasetIV2a.git} \\
\hline
\end{tabular}
\end{table}

\xmusection{超参数设置}{Hyperparameter Settings}
\label{sec:hyperparameters}

\xmusubsection{表格型 Q-Learning 超参数}{Tabular Q-Learning Hyperparameters}
\label{subsec:tabular_ql_params}

表格型 Q-Learning 方法（包括 Standard Q-Learning、Double Q-Learning、Dueling Q-Learning、Dueling Double Q-Learning）的超参数设置如表~\ref{tab:tabular_ql_params_state}--\ref{tab:tabular_ql_params_train} 所示。

\begin{table}[H]
\centering
\caption{表格型 Q-Learning 状态空间参数}
\label{tab:tabular_ql_params_state}
\begin{tabular}{ll}
\hline
\textbf{参数} & \textbf{值} \\
\hline
$t_{\text{start}}$ 范围 & [0.0, 3.0] 秒 \\
$t_{\text{end}}$ 范围 & [1.0, 4.0] 秒 \\
时间窗离散化步长 $\Delta t$ & 0.2 秒 \\
最小时间窗长度 & 0.5 秒 \\
状态空间大小 $|\mathcal{S}|$ & 136 \\
\hline
\end{tabular}
\end{table}

\begin{table}[H]
\centering
\caption{表格型 Q-Learning 学习参数}
\label{tab:tabular_ql_params_learn}
\begin{tabular}{ll}
\hline
\textbf{参数} & \textbf{值} \\
\hline
学习率 $\alpha$ & 0.2 \\
折扣因子 $\gamma$ & 0.95 \\
\hline
\end{tabular}
\end{table}

\begin{table}[H]
\centering
\caption{表格型 Q-Learning 探索策略参数}
\label{tab:tabular_ql_params_explore}
\begin{tabular}{ll}
\hline
\textbf{参数} & \textbf{值} \\
\hline
初始探索率 $\varepsilon_{\text{init}}$ & 0.9 \\
最小探索率 $\varepsilon_{\text{min}}$ & 0.05 \\
探索率衰减率 $\varepsilon_{\text{decay}}$ & 0.995 \\
\hline
\end{tabular}
\end{table}

\begin{table}[H]
\centering
\caption{表格型 Q-Learning 训练参数}
\label{tab:tabular_ql_params_train}
\begin{tabular}{ll}
\hline
\textbf{参数} & \textbf{值} \\
\hline
训练轮数 $n_{\text{episodes}}$ & 100 \\
每轮最大步数 & 40 \\
动作空间大小 $|\mathcal{A}|$ & 4 \\
\hline
\end{tabular}
\end{table}

\noindent\textbf{参数选择依据}：

（1）\textbf{状态空间离散化}：时间窗起始点 $t_{\text{start}}$ 以 0.2 秒步长在 [0.0, 3.0] 秒范围内离散化（16 个取值），结束点 $t_{\text{end}}$ 在 [1.0, 4.0] 秒范围内离散化（16 个取值），约束 $t_{\text{end}} - t_{\text{start}} \geq 0.5$ 秒，有效状态数量为 136 个。

（2）\textbf{学习率 $\alpha = 0.2$}：参考经典 Q-Learning 文献的常用设置，在收敛速度与稳定性之间取得平衡。

（3）\textbf{折扣因子 $\gamma = 0.95$}：使智能体在关注即时分类性能的同时，兼顾长期累积奖励。

（4）\textbf{探索策略}：初始探索率 $\varepsilon_{\text{init}} = 0.9$ 鼓励充分探索状态空间；衰减率 $\varepsilon_{\text{decay}} = 0.995$ 确保探索率缓慢下降；最小探索率 $\varepsilon_{\text{min}} = 0.05$ 避免陷入局部最优。

（5）\textbf{训练规模}：100 轮 × 40 步 = 4000 步的总探索次数，足以使 544 个状态 - 动作对（136 状态 × 4 动作）被充分访问（平均每对约 7.4 次）。

\xmusubsection{DQN 超参数}{DQN Hyperparameters}
\label{subsec:dqn_params}

DQN 基线方法的超参数设置如表~\ref{tab:dqn_params_state}--\ref{tab:dqn_params_replay} 所示。

\begin{table}[H]
\centering
\caption{DQN 状态空间参数}
\label{tab:dqn_params_state}
\begin{tabular}{ll}
\hline
\textbf{参数} & \textbf{值} \\
\hline
$t_{\text{start}}$ 范围 & [0.0, 3.0] 秒 \\
$t_{\text{end}}$ 范围 & [1.0, 4.0] 秒 \\
时间窗离散化步长 $\Delta t$ & 0.2 秒 \\
最小时间窗长度 & 0.5 秒 \\
\hline
\end{tabular}
\end{table}

\begin{table}[H]
\centering
\caption{DQN 学习参数}
\label{tab:dqn_params_learn}
\begin{tabular}{ll}
\hline
\textbf{参数} & \textbf{值} \\
\hline
学习率 $\eta$ & $1 \times 10^{-3}$ \\
折扣因子 $\gamma$ & 0.95 \\
\hline
\end{tabular}
\end{table}

\begin{table}[H]
\centering
\caption{DQN 探索策略参数}
\label{tab:dqn_params_explore}
\begin{tabular}{ll}
\hline
\textbf{参数} & \textbf{值} \\
\hline
初始探索率 $\varepsilon_{\text{init}}$ & 0.9 \\
最小探索率 $\varepsilon_{\text{min}}$ & 0.05 \\
探索率衰减率 $\varepsilon_{\text{decay}}$ & 0.995 \\
\hline
\end{tabular}
\end{table}

\begin{table}[H]
\centering
\caption{DQN 训练参数}
\label{tab:dqn_params_train}
\begin{tabular}{ll}
\hline
\textbf{参数} & \textbf{值} \\
\hline
训练轮数 $n_{\text{episodes}}$ & 100 \\
每轮最大步数 & 40 \\
\hline
\end{tabular}
\end{table}

\begin{table}[H]
\centering
\caption{DQN 经验回放参数}
\label{tab:dqn_params_replay}
\begin{tabular}{ll}
\hline
\textbf{参数} & \textbf{值} \\
\hline
缓冲区容量 & 10000 \\
批量大小 $\text{batch\_size}$ & 32 \\
目标网络更新频率 & 每 4 步 \\
软更新参数 $\tau$ & $1 \times 10^{-3}$ \\
\hline
\end{tabular}
\end{table}

\noindent\textbf{神经网络架构}：

DQN 使用两层全连接网络近似 Q 函数：

\textbf{输入层}：状态维度（时间窗参数编码）；

\textbf{隐藏层 1}：64 隐藏单元 + ReLU 激活；

\textbf{隐藏层 2}：32 隐藏单元 + ReLU 激活；

\textbf{输出层}：动作空间维度（4 维，对应 4 个动作的 Q 值）。

\xmusubsection{Dueling 架构参数}{Dueling Architecture Parameters}
\label{subsec:dueling_params}

Dueling Q-Learning 和 Dueling Double Q-Learning 采用 Dueling 网络架构，将 Q 函数分解为状态价值函数 $V(s)$ 和优势函数 $A(s,a)$：

\begin{equation}
Q(s, a) = V(s) + \left( A(s, a) - \frac{1}{|\mathcal{A}|} \sum_{a'} A(s, a') \right)
\end{equation}

Dueling 架构的学习参数如表~\ref{tab:dueling_params} 所示。

\begin{table}[H]
\centering
\caption{Dueling 架构学习参数}
\label{tab:dueling_params}
\begin{tabular}{ll}
\hline
\textbf{参数} & \textbf{值} \\
\hline
状态价值学习率 $\alpha_v$ & 0.05 \\
优势函数学习率 $\alpha_a$ & 0.1 \\
\hline
\end{tabular}
\end{table}

\noindent\textbf{说明}：状态价值函数和优势函数采用不同的学习率，以平衡两者对学习过程的贡献。

\xmusubsection{CSP-SVM 分类器超参数}{CSP-SVM Classifier Hyperparameters}
\label{subsec:csp_svm_params}

本研究使用 CSP-SVM 作为基础分类器计算奖励信号，其超参数设置如表~\ref{tab:csp_params}--\ref{tab:cv_params} 所示。为公平评估不同时间窗选择方法的性能，所有实验方法统一使用固定参数的 SVM 分类器，避免分类器参数变化对时间窗评估产生干扰。

\begin{table}[H]
\centering
\caption{CSP 特征提取参数}
\label{tab:csp_params}
\begin{tabular}{ll}
\hline
\textbf{参数} & \textbf{值} \\
\hline
CSP 分量数 & 4 \\
协方差矩阵正则化 & $1 \times 10^{-6}$ \\
特征类型 & 对数方差 \\
\hline
\end{tabular}
\end{table}

\begin{table}[H]
\centering
\caption{SVM 分类器参数}
\label{tab:svm_params}
\begin{tabular}{ll}
\hline
\textbf{参数} & \textbf{值} \\
\hline
核函数 & 线性核 \\
正则化参数 $C$ & 1.0 \\
随机种子 & 0 \\
\hline
\end{tabular}
\end{table}

\begin{table}[H]
\centering
\caption{交叉验证参数}
\label{tab:cv_params}
\begin{tabular}{ll}
\hline
\textbf{参数} & \textbf{值} \\
\hline
交叉验证折数 & 5 \\
\hline
\end{tabular}
\end{table}

\xmusection{源代码仓库}{Source Code Repository}
\label{sec:source_code}

本研究的全部源代码、实验配置及数据处理脚本均已开源，托管于 GitHub 平台。\textbf{GitHub 仓库地址}：\url{https://github.com/theoldman-lab/XMU_DKX.git}

\end{document}
